\section{Conclusions}

We introduce ACT-ViT, a fast and effective method for LLM Hallucination Detection (HD) that
relies solely on a model's internal representation of a single response. Unlike traditional probing,
ACT-ViT leverages the full Activation Tensor (AT) and, drawing on an analogy with images, it applies
techniques from modern vision architectures, enabling both cross-LLM training and generalization.
It is dramatically faster than multi-query methods (runtime $\approx 10^{-5}$ seconds), making it suitable for
real-time use. Across 15 LLM--dataset pairs, ACT-ViT consistently outperforms prior white-box
methods, supports zero-shot generalization to new datasets, and adapts efficiently to unseen LLMs.
Our work opens several directions for future research. While focused on HD, ACT-ViT can be
extended to tasks like data contamination and LLM-generated content detection. More advanced
architectures also warrant exploration, e.g., replacing per-LLM adapters with a shared module that
exploits permutation symmetries across ATs.

\paragraph{Limitations.} Due to the large size of the ATs, we begin their processing by conducting a simple
pooling operation to reduce their dimensionality. While effective in managing computational costs,
this step may discard potentially informative signals. Developing more sophisticated approaches to
handle AT size without sacrificing signal fidelity remains a promising future direction.

\paragraph{Acknowledgements.}
G.B. is supported by the Jacobs Qualcomm PhD Fellowship. F.F. conducted this work supported by
an Aly Kaufman Post-Doctoral Fellowship. H.M. is a Robert J. Shillman Fellow and is supported by
the Israel Science Foundation through a personal grant (ISF 264/23) and an equipment grant (ISF
532/23). F.F. is extremely grateful to the members of the ``Eva Project'', whose support he immensely
appreciates.
